\documentclass[11pt,a4paper]{article}
\usepackage[utf8]{inputenc}
\usepackage[T1]{fontenc}
\usepackage[spanish]{babel}
\usepackage{hyperref}
\usepackage{geometry}
\geometry{margin=2.5cm}
\usepackage{parskip}

\title{340 hacer el esfuerzo al máximo}
\author{Alejandro Valenzuela Rivera\\\small Universidad Autónoma de Nuevo León}
\date{\today}
\begin{document}
\maketitle

\begin{quote}
\textbf{Sesgo declarado.} Este trabajo se sostiene desde el sesgo declarado de su autor, y no finge neutralidad. Lo que afirma se puede tumbar, y ese es el punto: el falsificador está escrito, no escondido.
\end{quote}

\section*{Resumen}
La capacidad no se limita por prudencia: se lleva al máximo con criterio, y el costo se paga.

\section{Por qué se sostiene}
El que no usa todo lo que puede deja sin probar qué es: su tope es el de otro.

\section{De dónde viene}
La autoridad: 'algo que tenga que ver con Gurren Lagann, donde se trataba de hacer esfuerzo al máximo'.

\section{El caso}
Gurren Lagann: el esfuerzo al máximo, sin pedir permiso.

\section{Método}
\textquestiondown{}El tope lo puso el criterio, o el miedo de alguien?

\section{Límites, y qué lo tumba}
Un caso donde no llevar al máximo salve algo.

\section{Lo que queda abierto}
\textquestiondown{}Se llevó la capacidad a su tope?

\section*{Cómo se leyó esta tesis}
Esto se escribió en el repositorio man-made-empathy, del que este candado es una pieza: candados/06 del caso/340 hacer el esfuerzo al máximo.md. La máquina que lo compuso lo hizo desde este mismo archivo, y no agregó nada que no estuviera escrito. Si algo de esto es falso, se dice en el mismo sitio, y se corrige con fecha.

\end{document}
